\UTIResumenPagina
  {spanish}
  {Resumen Ejecutivo}
  {Las instituciones de educación superior (IES) del Ecuador gestionan datos personales y servicios académicos digitales, pero con frecuencia carecen de protocolos formales para detectar, gestionar y notificar incidentes de ciberseguridad. Esta carencia retrasa la respuesta, provoca la pérdida de evidencia y las expone a incumplir el plazo de 72 horas de la Ley Orgánica para el Fortalecimiento de la Ciberseguridad, cuando les es aplicable, y el término de cinco días para notificar a la Superintendencia de Protección de Datos Personales. Este trabajo desarrolla una metodología para la implementación de un protocolo de gestión y notificación de incidentes de ciberseguridad en las universidades e institutos superiores del país, siguiendo la investigación en ciencia del diseño. La metodología se estructura en cuatro etapas. Primero, un diagnóstico documental analiza la normativa vigente, los estudios previos y la información pública de una muestra de diez IES. Segundo, se diseñan siete fases de gestión (preparación, detección, clasificación, respuesta, notificación, cierre y lecciones aprendidas) que integran ISO/IEC 27035 y NIST SP 800-61 Rev.~3 con los plazos legales nacionales. Tercero, se elabora una matriz RACI de referencia, adaptable a instituciones de distinto tipo y tamaño. Cuarto, el protocolo se valida mediante el juicio de al menos cinco expertos, apoyado en un caso ilustrativo de vulneración de datos personales y evaluado con el coeficiente V de Aiken. Como resultados se obtienen el protocolo de referencia con cuatro plantillas, el informe de validación y recomendaciones para su adopción. Se concluye que articular los estándares internacionales con la normativa ecuatoriana en un flujo con responsables, tiempos y registros definidos ofrece a las IES una base común para fortalecer su capacidad de respuesta y cumplir los plazos legales.}
  {Descriptores}
  {Ciberseguridad, educación superior, gestión de incidentes, protección de datos personales.}
